% =========================================================
% Material for Chapter 6 (Limitations and Future Work) — negative and qualifying findings
% about the error-conditioned scoring method (4.4.2) and the preference pairs (4.4.3).
% Kept out of Section 5.2 on purpose. Numbers: result_5_2.ipynb, Sections 4-6 and 8-9, and
% Final__Preference_data_Statstic_Summary__train_only.ipynb (Monte-Carlo chi-square).
% Figure fig:model_composition (result_5_2/fig_5_12_model_composition.pdf) and Table tab:candidate_pool (text/table_candidate_pool.tex)
% are placed in Section 5.2. Paragraph headings are suggestions.
% =========================================================

    \paragraph{Statistical Strength of the Error--Move Association}
    The utilities of Section~\ref{subsec:error_conditioned_preference_scoring} describe tendencies in the 693 training dialogues rather than an established association. A chi-square test of independence on the contingency table of Table~\ref{tab:first_turn_error_move_cooccurrence} gives \(\chi^{2} = 25.94\) with 18 degrees of freedom (asymptotic \(p = 0.101\)); because 10 of the 28 expected counts are below 5 and three are below 1, the \(p\)-value was also estimated with a Monte-Carlo procedure that keeps both margins fixed (200,000 resamples), which gives \(p = 0.100\). The effect size is small, with a Cram\'{e}r's \(V\) of 0.112. The null hypothesis of independence between student error type and first tutor move can therefore not be rejected at the 5\,\% level, and the direction of the association is what the PPMI table encodes, not its statistical reliability.

    \paragraph{Support Behind the Largest Utilities}
    As Figure~\ref{fig:associations} shows, the three largest utilities of the table, all for \textit{Telling} (0.47, 1.67, and 1.18), rest on three, four, and one first tutor turns. A single additional or missing \textit{Telling} turn would change these values substantially, which is the known sensitivity of PMI to rare events noted in Section~\ref{subsec:error_conditioned_preference_scoring}. The same three error types are the ones for which the scoring prefers complete-solution responses over rechecking prompts (Table~\ref{tab:example_pairs}), so the least supported utilities are also the ones with the largest effect on which responses are chosen.

    \paragraph{An Error Type Without a Usable Signal}
    For \textit{misunderstanding of the question}, the largest error type with 6,966 pairs (28.7\,\% of the dataset), the only positive utility is that of \textit{Generic} (0.075). Generated candidates almost never realize this move: \textit{Generic} is the top label of 0.04\,\% of the candidates of this error type, and their mean \textit{Generic} posterior is 0.003. The scores of all its candidates lie between 0.0001 and 0.068, the median raw score margin of its pairs is 0.0001, and 98.6\,\% of its pairs have a raw margin below 0.001 (Figure~\ref{fig:margin_by_pair}). The three pairs of the dataset whose two responses have exactly the same score also belong to this error type, and 559 candidates in total share their score with another candidate of the same context, so that their side in a pair depends on the sort order rather than on the score. The chosen and rejected labels of this error type therefore carry almost no error-conditioned information, although its pairs form more than a quarter of the data used for preference optimization. The chosen responses of this error type also carry more \textit{Probing} and \textit{Telling} posterior mass than the rejected ones (means of 0.18 against 0.11 and 0.10 against 0.03; Figure~\ref{fig:posterior_shift}), although no utility rewards these moves; the small \textit{Generic} posterior co-varies with the posterior mass outside \textit{Focus}.

    \paragraph{A Candidate Pool Dominated by \textit{Focus}}
    The candidate pool is nearly uniform across error types and is dominated by one move: of the 48,839 candidates, 45,104 have \textit{Focus} as their top label, 2,221 \textit{Probing}, 1,498 \textit{Telling}, and 16 \textit{Generic}, and the same proportions hold within every error type (Table~\ref{tab:candidate_pool}). A context holds two \textit{Telling}-labelled candidates among about 70, and 307 of the 693 contexts hold none. Two consequences follow. First, for the three \textit{Telling}-associated error types the chosen responses are still \textit{Focus}-labelled responses in 82--93\,\% of the pairs, and 82--90\,\% of all pairs have the same top label on both sides; the preference signal of these error types is a difference in posterior mass (a mean \textit{Telling} posterior of 0.11--0.13 against 0.02) rather than a change of the realized move; in the contexts without a \textit{Telling}-labelled candidate, the scoring can only prefer the \textit{Focus} responses with the larger \textit{Telling} share. Second, for the \textit{Focus}-associated error types the chosen side consists almost entirely of \textit{Focus}-labelled responses (99.9--100\,\%) and the rejected side to 85--90\,\%, so most of these pairs separate \textit{Focus} responses by the confidence of the classifier rather than by the move. Both consequences originate outside the scoring method, in the Stage~2 prompt of Section~\ref{subsec:candidate_tutor_response_generation}, which asks for tutoring responses of at most two sentences, and in the classifier of Section~\ref{subsec:selected_classifier}, which assigns 41\,\% of the \textit{Probing} utterances of its test set to \textit{Focus}.

    \paragraph{Conditioning at the Level of Move Families}
    Re-ranking the candidates of every context with the utility column of each other error type (not shown; \path{result_5_2.ipynb}, Section~6) reveals that the seven utility columns produce only three distinct preference directions. The columns of \textit{numerical calculation error} and of \textit{none of the above} select the same responses, because both have a positive utility for \textit{Focus} only among the moves that the candidates realize, and the column of \textit{reached the correct solution but proceeded further} agrees with them on 74--78\,\%; the columns of \textit{missing or incorrect factual knowledge} and \textit{unit conversion error} agree on 99\,\%. The preference dataset therefore distinguishes \textit{Focus}-associated, \textit{Telling}-associated, and \textit{Generic}-only error types, but not the individual error types within a family. Differences between error types of the same family that a preference-optimized model may show cannot be attributed to the scoring method.

    \paragraph{Generation Model as a Confound}
    The six generation models differ in the moves they realize: \textit{Telling} is the top label of 12.5\,\% of the responses of \path{gpt-5.6-terra} but of 0--4\,\% of the responses of the other five models, and \textit{Probing} of 8.2\,\% of the responses of \path{mistral-nemo:12b}. The error-conditioned scoring consequently prefers different generation models for different error types (Figure~\ref{fig:model_composition} in Section~\ref{subsec:associations_in_pairs}). Responses of \path{gpt-5.6-terra} make up 27--29\,\% of the chosen responses of the three \textit{Telling}-associated error types but only 4--6\,\% of the chosen responses of the three \textit{Focus}-associated error types, against a pool share of 15--16\,\%; responses of \path{gemini-3.5-flash-lite} make up 4\,\% of the chosen responses for \textit{missing or extra quantity} and 27\,\% for \textit{reached the correct solution but proceeded further}, against 16\,\% of the pool. The chosen and rejected sides of a pair thus differ not only in the tutor move but also, systematically, in the generation model and its style (formatting, use of inline mathematics, length: mean response length ranges from 36 words for \path{mistral-nemo:12b} to 65 words for \path{llama3.2:1b}). A preference-optimized model can learn these stylistic differences instead of, or in addition to, the intended move.


    \paragraph{Incomparable Score Scales}
    The scores of different error types live on different scales, because each is bounded by the largest utility of its error type (0.075 for \textit{misunderstanding of the question}, 1.667 for \textit{missing or incorrect factual knowledge}). Margins were therefore normalized by that maximum before any comparison or selection across error types, but the normalization does not equalize the distributions: the mean normalized margin of the first pair ranges from 0.37 to 0.71 among the six error types with a usable signal (Figure~\ref{fig:margin_by_pair}), so the same selection ratio retains pairs of different separation in different error types.

    \paragraph{Association, Not Pedagogical Validity}
    Finally, a chosen response is one that realizes the move most associated with the error type in the training dialogues, not one that was judged pedagogically better. Table~\ref{tab:example_pairs} makes the consequence concrete: for \textit{numerical calculation error} the scoring prefers a prompt to recheck one step over the complete solution, whereas for \textit{missing or incorrect factual knowledge} it prefers the complete solution over the prompt. Whether either preference is the better tutoring decision was not evaluated, and the association it rests on is, as noted above, weak and thinly supported for the \textit{Telling} cases.
